%%%%%%%%%%%%%%%%%%%%%%%%%%%%%%%%%%%%%%%%%%%%%%%%%%%%%%%%%%%%%%%%%%%%%%%%%%%%%%%
%                       CARGA DE LA CLASE DE DOCUMENTO                        %
%                                                                             %
% Las opciones admisibles son:                                                %
%      12pt / 11pt            (tamaño del cuerpo de letra; no usar 10pt)      %
%                                                                             %
% catalan/spanish/english     (idioma principal del trabajo)                  %
%                                                                             % 
% french/italian/german...    (si necesitáis usar otro idioma adicional)      %
%                                                                             %
% listoffigures               (El documento incluye un Índice de figuras)     %
% listoftables                (El documento incluye un Índice de tablas)      %
% listofquadres               (El documento incluye un Índice de cuadros)     %
% listofalgorithms            (El documento incluye un Índice de algoritmos)  %
%                                                                             %
%%%%%%%%%%%%%%%%%%%%%%%%%%%%%%%%%%%%%%%%%%%%%%%%%%%%%%%%%%%%%%%%%%%%%%%%%%%%%%%

\documentclass[11pt,spanish,listoffigures,listoftables]{tfgetsinf}

%%%%%%%%%%%%%%%%%%%%%%%%%%%%%%%%%%%%%%%%%%%%%%%%%%%%%%%%%%%%%%%%%%%%%%%%%%%%%%%
%                     CODIFICACIÓN DEL ARCHIVO FUENTE                         %
%                                                                             %
%    Windows suele usar 'ansinew'                                             %
%    en Linux es posible que sea 'latin1' o 'latin9'                          %
%    Pero lo más recomendable es usar utf8 (unicode 8)                        %
%                                          (si vuestro editor lo permite)     % 
%%%%%%%%%%%%%%%%%%%%%%%%%%%%%%%%%%%%%%%%%%%%%%%%%%%%%%%%%%%%%%%%%%%%%%%%%%%%%%%

\usepackage[utf8]{inputenc} 


%%%%%%%%%%%%%%%%%%%%%%%%%%%%%%%%%%%%%%%%%%%%%%%%%%%%%%%%%%%%%%%%%%%%%%%%%%%%%%%
%                       OTROS PAQUETES Y DEFINICIONES                         %
%                                                                             %
%%%%%%%%%%%%%%%%%%%%%%%%%%%%%%%%%%%%%%%%%%%%%%%%%%%%%%%%%%%%%%%%%%%%%%%%%%%%%%%

\usepackage{glossaries}
\usepackage{textcomp}
\usepackage{booktabs}
\usepackage{float}
\usepackage{enumitem}
\usepackage{graphicx}
\usepackage{subcaption}
\usepackage{tabularx}
\usepackage{pgfplots}
\pgfplotsset{compat=1.18}

% Listings para código
\usepackage{listings}
\lstset{
  basicstyle=\ttfamily\small,
  frame=single,
  breaklines=true,
  postbreak=\mbox{\textcolor{gray}{$\hookrightarrow$}\space},
  columns=fullflexible,
  keepspaces=true,
  numbers=none
}

% Bibliografía
\usepackage[backend=biber,style=numeric,sorting=none]{biblatex}
\addbibresource{bibliografia.bib}

% Hyperref siempre al final
\usepackage{hyperref}
\hypersetup{
  colorlinks=true,
  linkcolor=black,
  urlcolor=cyan,
  citecolor=black
}
%%%%%%%%%%%%%%%%%%%%%%%%%%%%%%%%%%%%%%%%%%%%%%%%%%%%%%%%%%%%%%%%%%%%%%%%%%%%%%%
%                          DATOS DEL TRABAJO                                  %
%                                                                             %
% título alumno, titor y curso académico                                      %
%%%%%%%%%%%%%%%%%%%%%%%%%%%%%%%%%%%%%%%%%%%%%%%%%%%%%%%%%%%%%%%%%%%%%%%%%%%%%%%

\title{Entrenamiento de Redes Neuronales: MNIST y CIFAR-10}
\author{Gómez Corral, Miquel}
% \tutor{No}
% \curs{MUIARFID}

%%%%%%%%%%%%%%%%%%%%%%%%%%%%%%%%%%%%%%%%%%%%%%%%%%%%%%%%%%%%%%%%%%%%%%%%%%%%%%%
%                     PARAULES CLAU/PALABRAS CLAVE/KEY WORDS                  %
%                                                                             %
% Independentment de la llengua del treball, s'hi han d'incloure              %
% les paraules clau i el resum en els tres idiomes                            %
%%%%%%%%%%%%%%%%%%%%%%%%%%%%%%%%%%%%%%%%%%%%%%%%%%%%%%%%%%%%%%%%%%%%%%%%%%%%%%%

\keywords{
    Clasificación de imágenes; MNIST; CIFAR-10; Redes Neuronales; Deep Learning;
} % Paraules clau 
{
   Traducción de textos; 
} % Palabras clave
{
    Text clasification.
} % Key words


%%%%%%%%%%%%%%%%%%%%%%%%%%%%%%%%%%%%%%%%%%%%%%%%%%%%%%%%%%%%%%%%%%%%%%%%%%%%%%%
%                                                                             %
%                              INICI DEL DOCUMENT                             %
%                                                                             %
%%%%%%%%%%%%%%%%%%%%%%%%%%%%%%%%%%%%%%%%%%%%%%%%%%%%%%%%%%%%%%%%%%%%%%%%%%%%%%%

\begin{document} 


%%%%%%%%%%%%%%%%%%%%%%%%%%%%%%%%%%%%%%%%%%%%%%%%%%%%%%%%%%%%%%%%%%%%%%%%%%%%%%%
%              RESUMENES DEL TFG EN VALENCIA, CASTELLA I ANGLES               %
%%%%%%%%%%%%%%%%%%%%%%%%%%%%%%%%%%%%%%%%%%%%%%%%%%%%%%%%%%%%%%%%%%%%%%%%%%%%%%%


% \begin{abstract}[spanish]

% \end{abstract}

%%%%%%%%%%%%%%%%%%%%%%%%%%%%%%%%%%%%%%%%%%%%%%%%%%%%%%%%%%%%%%%%%%%%%%%%%%%%%%%
%                                                                             %
%                              CONTENIDO DEL TREBAJO                          %
%                                                                             %
%%%%%%%%%%%%%%%%%%%%%%%%%%%%%%%%%%%%%%%%%%%%%%%%%%%%%%%%%%%%%%%%%%%%%%%%%%%%%%%

%%%%%%%%%%%%%%%%%%%%%%%%%%%%%%%%%%%%%%%%%%%%%%%%%%%%%%%%%%%%%%%%%%%%%%%%%%%%%%%
%                                  ÍDNICE                                     %
%%%%%%%%%%%%%%%%%%%%%%%%%%%%%%%%%%%%%%%%%%%%%%%%%%%%%%%%%%%%%%%%%%%%%%%%%%%%%%%
% \clearpage
% \tableofcontents
% \listoffigures
% \listoftables


\mainmatter

\chapter{Introducción}
Este trabajo tiene como objetivo la creación y entrenamiento desde cero de modelos de redes neuronales para la resolución de dos tareas de clasificación de imágenes: MNIST y CIFAR-10. Se usarán diferentes técnicas y arquitecturas para cada uno de los datasets, y se evaluará el rendimiento de los modelos entrenados.

\subsubsection{Hardware utilizado}
Los experimentos se han realizado en mi máquina personal: GPU NVIDIA RTX 5070 (12 GB VRAM), procesador AMD Ryzen 9 9000X, 64 GB de RAM y Ubuntu 24.04 LTS.

Todos los modelos han podido ser entrenados y evaluados sin problemas. 

\section{Código entregado}
Se ha entregado un repositorio con todo el código necesario para reproducir los experimentos realizados en este trabajo. El repositorio tiene como punto de entrada el archivo auto-documentado \texttt{app/main.py}. Con él, se pueden entrenar los modelos para ambas tareas y ajustar fácilmente los hiperparámetros con argumentos.

El resto de la estructura del repositorio se encuentra en el archivo \texttt{README.md} del repositorio. Ahí se explica también como instalar las dependencias necesarias y cómo ejecutar el código.

Para lanzar los dos experimentos que han conseguido los mejores resultados, se pueden usar los siguientes comandos.
\begin{lstlisting}[language=Bash]
python ./main.py mnist -des "Best experiment" -b 1024
python ./main.py cifar -des "Best experiment" -b 256
\end{lstlisting}

Ajustar el batch size (-b) según la memoria disponible en la GPU. Con 12 GB de VRAM, estos valores funcionan correctamente.

%%%%%%%%%%%%%%%%%%%%%%%%%%%%%%%%%%%%%%%%%%%%%%%%%%%%%%%%%%%%%%%%%%%%%%%%%%%%%
%                        Notebooks base                                     %
%%%%%%%%%%%%%%%%%%%%%%%%%%%%%%%%%%%%%%%%%%%%%%%%%%%%%%%%%%%%%%%%%%%%%%%%%%%%%
\chapter{Datasets}
\section{Ejercicio 1: MNIST}
El dataset de MNIST es un conjunto de datos muy conocido en el ámbito del aprendizaje automático y la visión por computadora. Consiste en imágenes de dígitos manuscritos (del 0 al 9) en escala de grises, con un tamaño de 28x28 píxeles. El conjunto de datos contiene un total de 70,000 imágenes, divididas en 60,000 imágenes para entrenamiento y 10,000 imágenes para prueba.

La red utilizada para este ejercicio es una 'Multilayer Perceptron' (MLP) con 4 capas ocultas, una skip layer personalizada, que usa BatchNorm1d entre capas, Dropout del 45\% y ReLU como función de activación.

Para su entrenamiento, se ha utilizado el optimizador AdamW, OneCycleLR como scheduler y CrossEntropyLoss con \texttt{label smoothing} de 0.1 como función de pérdida. El modelo se ha entrenado durante 200 épocas con un tamaño de lote de 1024, y se ha objetivo una accuracy del $99.47\%$.

A continuación, se muestran los resultados que se han obtenido durante la implementación de las técnicas hasta llegar al objetivo.

La configuración base ha sido la siguiente: Adam, CrossEntropyLoss, Learning rate decay, data augmentation con RandomAffine (degrees=15, translate=(0.1, 0.1), scale=(0.9, 1.1), shear=10), arquitectura MLP con capas: [784, 1024, 1024, 512, 10]

\begin{table}[H]
\centering
\begin{tabular}{lrr}
    \toprule
    \textbf{Técnica} & \textbf{Accuracy \%} & \textbf{Error \%} \\
    \midrule
    Base (MLP 784-1024-1024-512) & 97.57 & 2.43 \\
    \midrule
    Data augmentation adaptativo & 98.36 & 1.64 \\
    Skip connections & 98.74 & 1.26 \\
    Red más densa (784-1536-1536-512-256) & 99.13 & 0.87 \\
    Skip connections adicionales & 99.20 & 0.80 \\
    Label smoothing y scheduler & 99.38 & 0.62 \\
    Ajuste final de hiperparámetros & 99.47 & 0.53 \\
    \bottomrule
\end{tabular}
\caption{Resultados de accuracy para MNIST de las distintas técnicas aplicadas.}
\label{tab:accuracy_mnist}
\end{table}

Al final, el uso de una red neuronal más densa con skip connections, junto con técnicas de regularización como label smoothing, data augmentation, y un ajuste fino de hiperparámetros, ha permitido alcanzar una accuracy del $99.47\%$ en el conjunto de prueba, lo que nos deja por debajo del $0.6\%$ de error.
\section{Ejercicio 2: CIFAR-10}
El dataset CIFAR-10 es otro conjunto de datos ampliamente utilizado. Consiste en 60,000 imágenes a color (3 canales, RGB) de 32x32 píxeles, divididas en 10 clases diferentes (avión, automóvil, pájaro, gato, ciervo, perro, rana, caballo, barco y camión). El conjunto de datos se divide en 50,000 imágenes para entrenamiento y 10,000 imágenes para prueba.

A diferencia del ejercicio anterior, la red utilizada sí es convolucional. En concreto, se compone de 3 bloques convolucionales de la clase 'ResNet' provista. Esta capas de 64, 128 y 256 canales, están seguidas de una capa fully connected al final. Se usa BatchNorm2d y ReLUs entre capas convolucionales, mientras que las fully conected tiene un BatchNorm1d. La primera FF también tiene un Dropout del 50\% y ReLU como función de activación. La última capa de la convolución se aplana y se conecta a una capa fully connected que produce las 10 salidas correspondientes a las clases del dataset. 

Para su entrenamiento, se ha vuelto a usar el optimizador AdamW y CrossEntropyLoss con \texttt{label smoothing} de 0.1 como función de pérdida, pero ahora con CosineAnnealingLR como scheduler. El modelo se ha entrenado durante 100 épocas con un tamaño de lote de 256, y se ha objetivo una accuracy del $93.44\%$.

La configuración base ha sido la siguiente: AdamW, CrossEntropyLoss (smooth 0.1), OneCycleLR, RandomHorizontalFlip, RandomAffine(degrees=15, translate=(0.1, 0.1), scale=(0.9, 1.1), shear=10), Canales convoluciones: [3, 64, 128, 256]

\begin{table}[H]
\centering
\begin{tabular}{lrr}
	\toprule
	\textbf{Técnica} & \textbf{Accuracy \%} & \textbf{Error \%} \\
	\midrule
	Base & 84.72 & 15.28 \\
	\midrule
	Rand augment con flip y Crop & 86.41 & 13.59 \\
	Resnet & 88.17 & 11.83 \\
	Dropout & 88.66 & 11.34 \\
	Ruido a las imágenes (no funciona) & 78.12 & 21.88 \\
	CosineAnnealingLR & 90.50 & 9.50 \\
	Mejores hiperparámetros y más Epochs & 91.81 & 8.19 \\
	CNN Más densa (3-32-64-128-256-512) & 93.44 & 6.56 \\
	\bottomrule
\end{tabular}
\caption{Resultados de accuracy para CIFAR-10 de las distintas técnicas aplicadas.}
\label{tab:accuracy_ej2}
\end{table}


Al final, el use de una red convolucional más densa, junto con las técnicas de data augmentation y un buen ajuste de hiperparámetros, ha permitido alcanzar la accuracy objetivo del $93.44\%$ en el conjunto de prueba, lo que nos deja por debajo del 8\% de error.

%%%%%%%%%%%%%%%%%%%%%%%%%%%%%%%%%%%%%%%%%%%%%%%%%%%%%%%%%%%%%%%%%%%%%%%%%%%%%%%
%                                                                             %
%                                BIBLIOGRAFIA                                 %
%                                                                             %
%%%%%%%%%%%%%%%%%%%%%%%%%%%%%%%%%%%%%%%%%%%%%%%%%%%%%%%%%%%%%%%%%%%%%%%%%%%%%%%
\cleardoublepage
\printbibliography

%%%%%%%%%%%%%%%%%%%%%%%%%%%%%%%%%%%%%%%%%%%%%%%%%%%%%%%%%%%%%%%%%%%%%%%%%%%%%%%
%                                                                             %
%                                 APÉNDICESS                                  %
%                                                                             %
%%%%%%%%%%%%%%%%%%%%%%%%%%%%%%%%%%%%%%%%%%%%%%%%%%%%%%%%%%%%%%%%%%%%%%%%%%%%%%%

\APPENDIX


%%%%%%%%%%%%%%%%%%%%%%%%%%%%%%%%%%%%%%%%%%%%%%%%%%%%%%%%%%%%%%%%%%%%%%%%%%%%%%
%                       EJEMPLOS DE CADA TIPO DE FACTURA                     %
%%%%%%%%%%%%%%%%%%%%%%%%%%%%%%%%%%%%%%%%%%%%%%%%%%%%%%%%%%%%%%%%%%%%%%%%%%%%%%

% \chapter{Apéndice correspondencia notebooks}
% \label{appendix:notebooks}
% ...

%%%%%%%%%%%%%%%%%%%%%%%%%%%%%%%%%%%%%%%%%%%%%%%%%%%%%%%%%%%%%%%%%%%%%%%%%%%%%%%
%                              FIN DEL DOCUMENTO                              %
%%%%%%%%%%%%%%%%%%%%%%%%%%%%%%%%%%%%%%%%%%%%%%%%%%%%%%%%%%%%%%%%%%%%%%%%%%%%%%%
\end{document}

